% GENERATED FILE. Edit canonical lessons, then run npm run generate:books.
% canonical-source-hash: fnv1a64:4fe6112fa70ed6bb
% canonical-lessons: HI-C373-mumkin, HI-C373-sambhav, HI-C373-sambhavna, HI-C373-cukna, HI-C373-upayogi

\chapter{Is It Possible?}
\label{ch:hi-is-it-possible}
\hlchaptermodality{373}

\begin{chapteropening}
\textbf{By the end of this chapter:} \emph{I can say whether something is possible, and on what condition.}

Complete the last lesson of chapter 373: I can say whether something is possible, and on what condition.
\end{chapteropening}

\section[\texorpdfstring{\hi{मुमकिन} (mumkin)}{मुमकिन (mumkin)}]{\hi{मुमकिन} (mumkin) — possible}

\label{lesson:HI-C373-mumkin}



\begin{quote}
Before the new one: say the Hindi for a sin, then the Hindi for to hold on to.
\end{quote}

\subsection*{You'll want to know: \hi{मुमकिन}}
\textbf{\hi{मुमकिन}} — \emph{mumkin} — \textquotedblleft{}possible\textquotedblright{}.

\begin{quote}
\hi{अगर} \hi{हम} \hi{जल्दी} \hi{निकलेंगे} \hi{तो} \hi{दोपहर} \hi{तक} \hi{पहुँचना} \hi{मुमकिन} \hi{है।}
\end{quote}

\begin{quote}
\emph{agar ham jaldī nikleṅge to dopahar tak pahũcnā mumkin hai.}
\end{quote}

\begin{quote}
If we leave early, we can get there by noon.
\end{quote}

\begin{cousinweb}
\textbf{\hi{मुमकिन}} (\emph{mumkin}) is usually traced to Arabic \emph{mumkin}, \textquotedblleft{}possible\textquotedblright{}, which reached Hindi through Persian.
\end{cousinweb}

\begin{grammarlens}[title={What you've built}]
One more word for saying what happens if.
\end{grammarlens}

\subsection*{Your turn}
\begin{itemize}
\raggedright
  \item \emph{Say it:} \emph{mumkin}
  \item \emph{Say it:} \emph{mumkin}, once more
  \item \emph{Say it:} \emph{agar ham jaldī nikleṅge to dopahar tak pahũcnā mumkin hai.} — If we leave early, we can get there by noon.
  \item \emph{Recall:} say the Hindi for a sin, then the Hindi for to hold on to, then say \emph{mumkin} again
\end{itemize}

\begin{tcolorbox}[breakable,colback=teal!4,colframe=teal!35!black,title={Before you move on}]
What does \hi{मुमकिन} mean? (\textquotedblleft{}Possible\textquotedblright{}.) Say it once more.
\end{tcolorbox}
\section[\texorpdfstring{\hi{संभव} (sambhav)}{संभव (sambhav)}]{\hi{संभव} (sambhav) — possible (formal)}

\label{lesson:HI-C373-sambhav}



\begin{quote}
Before the new one: say the Hindi for to hold on to, then the Hindi for possible.
\end{quote}

\subsection*{You'll want to know: \hi{संभव}}
\textbf{\hi{संभव}} — \emph{sambhav} — \textquotedblleft{}possible (formal)\textquotedblright{}.

\begin{quote}
\hi{अगर} \hi{सब} \hi{मदद} \hi{करेंगे} \hi{तो} \hi{यह} \hi{काम} \hi{संभव} \hi{है।}
\end{quote}

\begin{quote}
\emph{agar sab madad kareṅge to yah kām sambhav hai.}
\end{quote}

\begin{quote}
If everyone helps, this job can be done.
\end{quote}

\begin{cousinweb}
\textbf{\hi{संभव}} (\emph{sambhav}) is a tatsama, borrowed from Sanskrit \emph{sambhava}, \textquotedblleft{}coming into being\textquotedblright{} (Monier-Williams). Hindi uses it as an adjective: what can come about.
\end{cousinweb}

\begin{grammarlens}[title={What you've built}]
One more word for saying what happens if.
\end{grammarlens}

\subsection*{Your turn}
\begin{itemize}
\raggedright
  \item \emph{Say it:} \emph{sambhav}
  \item \emph{Say it:} \emph{sambhav}, once more
  \item \emph{Say it:} \emph{agar sab madad kareṅge to yah kām sambhav hai.} — If everyone helps, this job can be done.
  \item \emph{Recall:} say the Hindi for to hold on to, then the Hindi for possible, then say \emph{sambhav} again
\end{itemize}

\begin{tcolorbox}[breakable,colback=teal!4,colframe=teal!35!black,title={Before you move on}]
What does \hi{संभव} mean? (\textquotedblleft{}Possible (formal)\textquotedblright{}.) Say it once more.
\end{tcolorbox}
\section[\texorpdfstring{\hi{संभावना} (sambhāvnā)}{संभावना (sambhāvnā)}]{\hi{संभावना} (sambhāvnā) — a likelihood, a possibility}

\label{lesson:HI-C373-sambhavna}



\begin{quote}
Before the new one: say the Hindi for possible, then the Hindi for possible (formal).
\end{quote}

\subsection*{You'll want to know: \hi{संभावना}}
\textbf{\hi{संभावना}} — \emph{sambhāvnā} — \textquotedblleft{}a likelihood, a possibility\textquotedblright{}.

\begin{quote}
\hi{अगर} \hi{बादल} \hi{रहेंगे} \hi{तो} \hi{बारिश} \hi{की} \hi{संभावना} \hi{है।}
\end{quote}

\begin{quote}
\emph{agar bādal raheṅge to bāriś kī sambhāvnā hai.}
\end{quote}

\begin{quote}
If the clouds stay, there is a chance of rain.
\end{quote}

Feminine.

\begin{cousinweb}
\textbf{\hi{संभावना}} (\emph{sambhāvnā}) is a tatsama, borrowed from Sanskrit \emph{sambhāvanā}, \textquotedblleft{}bringing about; supposing\textquotedblright{} (Monier-Williams).
\end{cousinweb}

\begin{grammarlens}[title={What you've built}]
One more word for saying what happens if.
\end{grammarlens}

\subsection*{Your turn}
\begin{itemize}
\raggedright
  \item \emph{Say it:} \emph{sambhāvnā}
  \item \emph{Say it:} \emph{sambhāvnā}, once more
  \item \emph{Say it:} \emph{agar bādal raheṅge to bāriś kī sambhāvnā hai.} — If the clouds stay, there is a chance of rain.
  \item \emph{Recall:} say the Hindi for possible, then the Hindi for possible (formal), then say \emph{sambhāvnā} again
\end{itemize}

\begin{tcolorbox}[breakable,colback=teal!4,colframe=teal!35!black,title={Before you move on}]
What does \hi{संभावना} mean? (\textquotedblleft{}A likelihood, a possibility\textquotedblright{}.) Say it once more.
\end{tcolorbox}
\section[\texorpdfstring{\hi{चूकना} (cūknā)}{चूकना (cūknā)}]{\hi{चूकना} (cūknā) — to miss (a chance, a mark)}

\label{lesson:HI-C373-cukna}



\begin{quote}
Before the new one: say the Hindi for possible (formal), then the Hindi for a likelihood.
\end{quote}

\subsection*{You'll want to know: \hi{चूकना}}
\textbf{\hi{चूकना}} — \emph{cūknā} — \textquotedblleft{}to miss (a chance, a mark)\textquotedblright{}.

\begin{quote}
\hi{अगर} \hi{तुम} \hi{तेज़} \hi{नहीं} \hi{दौड़ोगे} \hi{तो} \hi{गेंद} \hi{पकड़ने} \hi{से} \hi{चूक} \hi{जाओगे।}
\end{quote}

\begin{quote}
\emph{agar tum tez nahīṃ dauṛoge to gẽd pakaṛne se cūk jāoge.}
\end{quote}

\begin{quote}
If you don't run fast, you'll fail to catch the ball.
\end{quote}

Past (he …): \hi{वह} \hi{चूक} \hi{गया}, \emph{vah cūk gayā}.

\begin{cousinweb}
\textbf{\hi{चूकना}} (\emph{cūknā}) is inherited from \emph{*cukk}, \textquotedblleft{}fall short\textquotedblright{}, a form CDIAL reconstructs from its descendants (CDIAL 4848).
\end{cousinweb}

\begin{grammarlens}[title={What you've built}]
One more word for saying what happens if.
\end{grammarlens}

\subsection*{Your turn}
\begin{itemize}
\raggedright
  \item \emph{Say it:} \emph{cūknā}
  \item \emph{Say it:} \emph{cūknā}, once more
  \item \emph{Say it:} \emph{agar tum tez nahīṃ dauṛoge to gẽd pakaṛne se cūk jāoge.} — If you don't run fast, you'll fail to catch the ball.
  \item \emph{Recall:} say the Hindi for possible (formal), then the Hindi for a likelihood, then say \emph{cūknā} again
\end{itemize}

\begin{tcolorbox}[breakable,colback=teal!4,colframe=teal!35!black,title={Before you move on}]
What does \hi{चूकना} mean? (\textquotedblleft{}To miss (a chance, a mark)\textquotedblright{}.) Say it once more.
\end{tcolorbox}
\section[\texorpdfstring{\hi{उपयोगी} (upayogī)}{उपयोगी (upayogī)}]{\hi{उपयोगी} (upayogī) — useful}

\label{lesson:HI-C373-upayogi}



\begin{quote}
Before the new one: say the Hindi for a likelihood, then the Hindi for to miss.
\end{quote}

\subsection*{You'll want to know: \hi{उपयोगी}}
\textbf{\hi{उपयोगी}} — \emph{upayogī} — \textquotedblleft{}useful\textquotedblright{}.

\begin{quote}
\hi{अगर} \hi{तुम} \hi{यह} \hi{किताब} \hi{पढ़ोगे} \hi{तो} \hi{यह} \hi{बहुत} \hi{उपयोगी} \hi{होगी।}
\end{quote}

\begin{quote}
\emph{agar tum yah kitāb paṛhoge to yah bahut upayogī hogī.}
\end{quote}

\begin{quote}
If you read this book, you will find it very useful.
\end{quote}

\begin{cousinweb}
\textbf{\hi{उपयोगी}} (\emph{upayogī}) is a tatsama, borrowed from Sanskrit \emph{upayogin}, \textquotedblleft{}serviceable, useful\textquotedblright{} (Monier-Williams).
\end{cousinweb}

\begin{grammarlens}[title={What you've built}]
One more word for saying what happens if.
\end{grammarlens}

\subsection*{Your turn}
\begin{itemize}
\raggedright
  \item \emph{Say it:} \emph{upayogī}
  \item \emph{Say it:} \emph{upayogī}, once more
  \item \emph{Say it:} \emph{agar tum yah kitāb paṛhoge to yah bahut upayogī hogī.} — If you read this book, you will find it very useful.
  \item \emph{Recall:} say the Hindi for a likelihood, then the Hindi for to miss, then say \emph{upayogī} again
\end{itemize}

\begin{tcolorbox}[breakable,colback=teal!4,colframe=teal!35!black,title={Before you move on}]
What does \hi{उपयोगी} mean? (\textquotedblleft{}Useful\textquotedblright{}.) Say it once more.
\end{tcolorbox}
